% Generated draft: requires mathematical and source review.
\documentclass[11pt]{article}
\usepackage[margin=25mm]{geometry}
\usepackage{fontspec}
\setmainfont{Latin Modern Roman}
\usepackage{amsmath,amssymb,mathtools,unicode-math}
\setmathfont{Latin Modern Math}
\usepackage{enumitem,needspace,xurl,hyperref}
\setlength{\parindent}{0pt}
\setlength{\parskip}{5pt}
\setlength{\emergencystretch}{3em}
\newcommand{\N}{\mathbb{N}}
\newcommand{\Z}{\mathbb{Z}}
\newcommand{\Q}{\mathbb{Q}}
\newcommand{\R}{\mathbb{R}}
\newcommand{\C}{\mathbb{C}}
\newcommand{\F}{\mathbb{F}}
\newcommand{\A}{\mathbb{A}}
\newcommand{\PP}{\mathbb{P}}
\newcommand{\E}{\mathbb{E}}
\newcommand{\eps}{\varepsilon}
\newcommand{\dd}{\,\mathrm{d}}
\begin{document}
\section*{UnsolvedMath ID: 800801}
\textbf{Source problem code:} catalogue entry
\subsection{Definitions and mathematical statement}
The supplied excerpt does not specify the class of objects being quantified over or define the symbols \(f_0\) and \(f_d\). Consequently, it does not determine a complete mathematical conjecture.

Under the usual simplicial face-count interpretation, for a finite abstract simplicial complex \(K\), define
\[
f_i(K)=\#\{\sigma\in K:|\sigma|=i+1\},\qquad i\geq 0.
\]
Thus \(f_0(K)\) counts vertices and \(f_d(K)\) counts \(d\)-dimensional faces. This interpretation is provisional; the excerpt does not explicitly identify the objects as simplicial complexes.

Writing \(\mathcal{C}_d\) for the intended, unspecified class of objects, the supplied inequality has the following quantified template:
\[
\forall d\in\Z_{\geq 1}\quad \exists C_d\in\R_{>0}\quad
\forall K\in\mathcal{C}_d,\qquad
f_d(K)\leq C_d\bigl(f_0(K)\bigr)^d.
\]
Here \(C_d\) depends only on \(d\), not on \(K\). The class \(\mathcal{C}_d\), its geometric or topological restrictions, and any intended ambient dimension remain unresolved. The template must therefore not be read as a conjecture about all finite simplicial complexes.

The source presents this as another possible generalization of the planar-graph inequality \(e\leq 3v\), where \(e\) and \(v\) denote the numbers of edges and vertices. It reports that work of Kalai and Sarkaria suggests a more precise formulation that would yield the best possible constant in each dimension, but that formulation is not included in the supplied record.
\subsection{Short English statement}
For every integer \(d\geq 1\), the proposed bound asks for a constant depending only on \(d\) such that the number of \(d\)-dimensional faces is at most that constant times the \(d\)th power of the number of vertices. The intended class of objects and its hypotheses are missing, so this remains an incomplete formulation requiring source review.
\subsection{Sources}
\begin{itemize}
\item Local UnsolvedMath registry: permanent ID \(800801\), source problem code \texttt{catalogue entry}, title \emph{Branko Grunbaum in Many Dimensions --- Conjecture 8}. The record attributes the problem to Kahle, with proposed year 2019, and labels it open. This status is supplied metadata, not an independently verified mathematical assessment.
\item Supplied source list: Kahle, \emph{Branko Grunbaum in Many Dimensions: Open Problems} (2019), Conjecture 8. Supplied URL: \url{https://arxiv.org/abs/1901.08622}. The record reports access on 2026-10-16 and extraction by \texttt{arxiv\_latex}; these are supplied provenance details, not a claim of independent source access or verification.
\item The excerpt cites Dey [Dey93] for earlier discussion of various formulations, and Kalai [Kalai91] and Sarkaria [Sarkaria92] for work suggesting a precise formulation and optimal constants. Full bibliographic details and the referenced formulation are absent from the supplied record.
\item The excerpt reports that Gr\"unbaum described the question as still open in a 1970 article. It also preserves Kahle's qualification that he had not found an earlier reference. No bibliographic details for that article are supplied, and this historical report has not been independently verified.
\end{itemize}
\end{document}
